```latex
% ============================================================
% TITLE PAGE — SGCV-IA — Requirements Engineering (Springer)
%
% Archivo independiente con información de autoría y declaraciones.
% Mantener separado del manuscrito anonimizado cuando el proceso de
% revisión de la revista así lo requiera.
% ============================================================

\documentclass[12pt]{article}

\usepackage[margin=2.5cm]{geometry}
\usepackage[T1]{fontenc}
\usepackage[utf8]{inputenc}
\usepackage{hyperref}

\begin{document}

% ============================================================
% TITLE
% ============================================================

\begin{center}
{\Large\bfseries
Explainability as a non-functional requirement in an AI-assisted
veterinary clinical system: a case study
}
\end{center}

\vspace{1em}

% ============================================================
% AUTHORS
% ============================================================

\section*{Authors}

\begin{enumerate}

  \item Alberto Jeanpool Marcillo Ponce$^{1,\ast}$\\
  ORCID: \href{https://orcid.org/0009-0003-7886-3957}
  {0009-0003-7886-3957}\\
  \texttt{amarcillop@uteq.edu.ec}

  \item Robyn Willian Amagua Sacón$^{1}$\\
  \texttt{ramaguas@uteq.edu.ec}

  \item Anthony Alfredo Vera Gómez$^{1}$\\
  \texttt{averag10@uteq.edu.ec}

  \item Jhon Alexander Mesías Quijije$^{1}$\\
  \texttt{jmesiasq@uteq.edu.ec}

  \item Carlos Daniel Barrionuevo Fuentes$^{1}$\\
  \texttt{cbarrionuevof@uteq.edu.ec}

  \item Gleiston Ciceron Guerrero Ulloa$^{1}$\\
  \texttt{gguerrero@uteq.edu.ec}

\end{enumerate}

\noindent
$^{1}$Facultad de Ciencias de la Computación,
Universidad Técnica Estatal de Quevedo (UTEQ),
Quevedo, Ecuador.

\vspace{0.5em}

\noindent
$^{\ast}$Corresponding author:
Alberto Jeanpool Marcillo Ponce,
\texttt{amarcillop@ute.edu.ec}

\vspace{0.5em}

\noindent
\textit{Note: ORCID identifiers for authors other than Alberto Jeanpool
Marcillo Ponce have not yet been provided.}

\vspace{1em}

% ============================================================
% ABSTRACT
% ============================================================

\section*{Abstract}

The final abstract must be identical to the abstract included in the
final manuscript. It will be inserted after completion and verification
of the terminal empirical analysis, since the Results and Conclusions
components depend on the experimental evidence.

% ============================================================
% KEYWORDS
% ============================================================

\section*{Keywords}

Explainability; Non-functional requirements; AI-assisted diagnosis;
Veterinary informatics; User studies

\vspace{1.5em}

% ============================================================
% DECLARATIONS
% ============================================================

\section*{Declarations}

% ------------------------------------------------------------
% FUNDING
% ------------------------------------------------------------

\subsection*{Funding}

\textit{Pending confirmation before submission.}
The authors must state whether the study received financial or material
support from UTEQ, Veterinaria Macay, or any other organization.
If no funding was received, the final statement should indicate that
the authors received no specific funding for this work.

% ------------------------------------------------------------
% COMPETING INTERESTS
% ------------------------------------------------------------

\subsection*{Competing Interests}

\textit{Pending confirmation before submission.}
The authors must declare any relevant financial or non-financial
interests, including any relevant relationship between the project
team and the organization associated with the case study.

% ------------------------------------------------------------
% ETHICS
% ------------------------------------------------------------

\subsection*{Ethics approval and consent to participate}

Participation in the requirements-elicitation activities was voluntary
and based on written informed consent. Participant information is
handled using anonymized identifiers, and identifying materials are
maintained separately in the restricted evidence area in accordance
with the project's data-protection procedures and Ecuador's
\textit{Ley Orgánica de Protección de Datos Personales}
(LOPDP).

\par
\textit{Pending confirmation before submission:}
the final manuscript must state whether formal approval was granted by
an institutional ethics committee or equivalent body and, if
applicable, provide the official approval identifier. No approval
number is reported here because none has been provided.

% ------------------------------------------------------------
% CONSENT FOR PUBLICATION
% ------------------------------------------------------------

\subsection*{Consent for publication}

\textit{Pending verification before submission.}
The final statement must reflect the exact scope of the signed consent
documents regarding the use and publication of anonymized participant
information.

% ------------------------------------------------------------
% DATA / MATERIALS / CODE
% ------------------------------------------------------------

\subsection*{Data, Materials and/or Code availability}

The research materials supporting the requirements-engineering phase
include anonymized interview transcripts, thematic-coding materials,
requirements artefacts, traceability information, and related project
documentation. Potentially identifying materials, including original
consent documents and multimedia evidence, are maintained separately
in the restricted evidence area and are not intended for public
dissemination.

\par
The terminal replication package is intended to include the
anonymized research materials and reproducible analysis artefacts
required by the project. The persistent Zenodo DOI and the
corresponding OSF registration identifier will be added once those
resources have been finalized. No DOI or OSF identifier is reported
until it has been verified.

% ------------------------------------------------------------
% AUTHOR CONTRIBUTIONS
% ------------------------------------------------------------

\subsection*{Author contributions}

\textit{Pending confirmation by the authors before submission.}
The final contribution statement will describe the actual contribution
of each author, preferably following the CRediT taxonomy.

% ------------------------------------------------------------
% AI-ASSISTED TECHNOLOGIES
% ------------------------------------------------------------

\subsection*{Use of AI-assisted technologies}

During preparation of the manuscript, ChatGPT (OpenAI, GPT-5.6 Luna)
was used to assist with drafting, restructuring, language refinement,
and consistency checking of manuscript text based on information and
research artefacts supplied by the authors.

The AI-assisted tool was not used to fabricate empirical observations,
participant responses, statistical results, tables, figures, effect
sizes, confidence intervals, or scientific conclusions. Quantitative
and qualitative findings included in the final manuscript must be
derived from the project's actual evidence and reproducible analysis
procedures.

% ------------------------------------------------------------
% ACKNOWLEDGEMENTS
% ------------------------------------------------------------

\subsection*{Acknowledgements}

The authors thank the veterinary professionals, administrative
personnel, and other participants who contributed to the
requirements-elicitation and validation activities of the SGCV-IA
project. Participant identities are not disclosed in this manuscript
in accordance with the project's anonymization and data-protection
procedures.

\end{document}
```
